% ============================================================
% SECTION: LITERATURE SURVEY
% ============================================================

\section{Literature Survey}

Machine learning techniques have been widely adopted in financial domains such as credit risk assessment, real estate valuation, and stock price prediction due to their ability to process large and complex datasets efficiently and uncover hidden patterns. These approaches have demonstrated strong predictive performance compared to traditional statistical methods, enabling faster and more data-driven financial decision-making. However, the lack of interpretability in many advanced models has limited their practical adoption in high-stakes financial environments, where accountability and transparency are critical. 

Xu et al. \cite{xu2023} proposed an explainable framework for loan default prediction by integrating machine learning models with feature attribution techniques. Their study showed that incorporating explainability improves transparency and user trust while maintaining competitive predictive accuracy. Rashid and Naeem \cite{rashid2021} provided a comprehensive analytical review of Explainable Artificial Intelligence (XAI) methods, categorizing them into model-specific and model-agnostic approaches, and emphasized their growing importance in regulated domains such as finance, healthcare, and governance. 

Tripathi et al. \cite{tripathi2021} conducted an experimental evaluation of multiple machine learning algorithms for credit score classification. Their results indicated that ensemble-based models outperform linear classifiers in terms of accuracy and robustness; however, these models often suffer from reduced transparency. Similarly, Coşer et al. \cite{coser2019} analyzed predictive models for loan default risk assessment and highlighted the trade-off between predictive performance and interpretability, stressing the need for explainable solutions in financial decision-support systems. 

In the real estate domain, Maha and Valluri \cite{maha2021} explored deep learning and natural language processing techniques for property-related analytics, demonstrating the potential of machine learning in automating real estate insights. While their approach achieved effective results, it did not integrate explainability mechanisms, which limits its applicability in real-world financial decision-making where justification of predictions is essential. 

Stock price prediction has also been extensively studied using neural networks, regression models, and hybrid approaches. Although these methods can achieve reasonable forecasting accuracy, they are typically treated as black-box systems, providing limited insight into the influence of market variables and historical trends on predictions. This opacity restricts their acceptance in practical investment and risk analysis scenarios. 

To address these challenges, explainability techniques have gained significant attention in recent research. Lundberg and Lee \cite{lundberg2017} introduced SHAP, a game-theoretic framework that offers consistent local and global explanations by quantifying feature contributions. Ribeiro et al. \cite{ribeiro2016} proposed LIME, a model-agnostic technique that explains individual predictions by locally approximating complex models. Molnar \cite{molnar2020} and Guidotti et al. \cite{guidotti2018} further emphasized the importance of interpretability in machine learning systems, particularly in domains governed by ethical, legal, and regulatory constraints. 

Despite these advancements, most existing studies focus on single-task financial applications and lack an integrated, interactive platform that supports multiple financial prediction tasks with real-time explainability. This paper addresses this gap by proposing a unified explainable AI--based financial decision-support system that simultaneously supports loan approval prediction, property price estimation, and stock price forecasting through a common web-based interface.
